\section{Theoretical Analysis}

In this section, we provide theoretical analysis of the sequential convexification algorithm, including convergence properties and tube-excess bounds.

\subsection{Problem Formulation}

Recall the optimization problem:
\begin{align}
\min_{\phi \in \mathbb{R}^{3M}} \quad & \frac{1}{2} \phi^T H \phi \\
\text{s.t.} \quad & \|r_i(\phi_i)\|_2 \leq \epsilon_i, \quad i = 0, \ldots, M-1 \nonumber
\end{align}

where $r_i(\phi_i) = \log(R_{\text{meas},i}^T \exp(\phi_i^\wedge))$ and $H$ is the smoothness Hessian.

\subsection{Sequential Convexification Convergence}

At each outer iteration $k$, we solve a convex subproblem by linearizing the constraints around the current iterate $\phi^k$:

\begin{align}
\min_{\delta \in \mathbb{R}^{3M}} \quad & \frac{1}{2} (\phi^k + \delta)^T H (\phi^k + \delta) \\
\text{s.t.} \quad & \|r_i + J_i \delta_i\|_2 \leq \epsilon_i, \quad i = 0, \ldots, M-1 \nonumber \\
& \|\delta_i\|_2 \leq \Delta, \quad i = 0, \ldots, M-1 \nonumber
\end{align}

where $r_i = r_i(\phi_i^k)$ and $J_i = \frac{\partial r_i}{\partial \phi_i}\big|_{\phi_i^k}$.

\begin{lemma}[Lipschitz Continuity of Residual Jacobian]
The residual Jacobian $J_r(\phi) = J_r^{-1}(\log(R_{\text{meas}}^T \exp(\phi^\wedge))) J_r(\phi)$ is Lipschitz continuous on compact sets. Specifically, for any $\phi_1, \phi_2$ with $\|\phi_1\|, \|\phi_2\| \leq \pi$:
\begin{equation}
\|J_r(\phi_1) - J_r(\phi_2)\| \leq L_J \|\phi_1 - \phi_2\|
\end{equation}
where $L_J = \frac{1}{2}$ is the Lipschitz constant.
\end{lemma}

\begin{proof}
The right Jacobian $J_r(\phi)$ has the form:
\begin{equation}
J_r(\phi) = I - \frac{1-\cos\|\phi\|}{\|\phi\|^2} \phi^\wedge + \frac{\|\phi\| - \sin\|\phi\|}{\|\phi\|^3} (\phi^\wedge)^2
\end{equation}

The derivatives of the coefficients $\alpha(\|\phi\|) = \frac{1-\cos\|\phi\|}{\|\phi\|^2}$ and $\beta(\|\phi\|) = \frac{\|\phi\| - \sin\|\phi\|}{\|\phi\|^3}$ are bounded for $\|\phi\| \leq \pi$, giving the Lipschitz constant $L_J = 1/2$.
\end{proof}

\begin{theorem}[Local Convergence of Sequential Convexification]
Assume:
\begin{enumerate}
\item The initial guess $\phi^0$ is sufficiently close to a local minimum $\phi^*$.
\item The trust region radius $\Delta$ satisfies $\Delta \leq \min(\Delta_{\max}, 1/(2L_J))$.
\item The linearized subproblems satisfy the Linear Independence Constraint Qualification (LICQ).
\end{enumerate}

Then the sequence $\{\phi^k\}$ generated by the sequential convexification algorithm converges to $\phi^*$ with linear rate:
\begin{equation}
\|\phi^{k+1} - \phi^*\| \leq \rho \|\phi^k - \phi^*\|
\end{equation}
where $\rho = 1 - \frac{\eta}{\kappa}$ depends on the acceptance threshold $\eta$ and condition number $\kappa$ of the Hessian $H$.
\end{theorem}

\begin{proof}[Proof Sketch]
The proof follows standard SQP convergence theory:
\begin{enumerate}
\item By Lemma 1, the linearization error is $O(\|\delta\|^2)$.
\item The trust region ensures $\|\delta\| \leq \Delta$.
\item For accepted steps ($\rho_k \geq \eta_{\min}$), the actual decrease is at least a fraction of the predicted decrease.
\item Near the solution, all steps are accepted and convergence is linear.
\end{enumerate}
\end{proof}

\subsection{Constraint Violation Bounds}

While the linearized constraints are satisfied exactly in each subproblem, the actual nonlinear constraints may be violated due to linearization error.

\begin{lemma}[Linearization Error Bound]
For the tube constraint $c_i(\phi) = \|r_i(\phi)\|_2 - \epsilon_i$, the linearization error is bounded by:
\begin{equation}
|c_i(\phi + \delta) - \tilde{c}_i(\delta)| \leq \frac{L_J}{2} \|\delta\|^2
\end{equation}
where $\tilde{c}_i(\delta) = \|r_i + J_i \delta\|_2 - \epsilon_i$ is the linearized constraint.
\end{lemma}

\begin{proof}
By Taylor's theorem with remainder and Lemma 1 (Lipschitz Jacobian):
\begin{align}
r_i(\phi + \delta) &= r_i(\phi) + J_i \delta + O(\|\delta\|^2) \\
\|r_i(\phi + \delta)\| &= \|r_i + J_i \delta\| + O(\|\delta\|^2)
\end{align}
The remainder term is bounded by $\frac{L_J}{2}\|\delta\|^2$.
\end{proof}

\begin{theorem}[Feasibility Guarantee]
If the trust region radius satisfies:
\begin{equation}
\Delta \leq \sqrt{\frac{2\epsilon_{\min}}{L_J}}
\end{equation}
where $\epsilon_{\min} = \min_i \epsilon_i$, then the sequential convexification algorithm maintains feasibility up to tolerance:
\begin{equation}
\|r_i(\phi^k)\|_2 \leq \epsilon_i + \frac{L_J \Delta^2}{2} \quad \forall i, k
\end{equation}
\end{theorem}

\begin{proof}
At each iteration, the linearized constraint is satisfied: $\|r_i + J_i \delta_i\|_2 \leq \epsilon_i$.
By Lemma 2, the actual tube excess is bounded by the linearization error:
\begin{equation}
\|r_i(\phi + \delta)\|_2 \leq \|r_i + J_i \delta\|_2 + \frac{L_J}{2}\|\delta\|^2 \leq \epsilon_i + \frac{L_J \Delta^2}{2}
\end{equation}
\end{proof}

\subsection{Practical Implications}

The theoretical results provide practical guidance:

\textbf{Trust Region Selection:}
Theorem 2 suggests choosing $\Delta$ such that $L_J \Delta^2 / 2 \ll \epsilon_{\min}$. For typical values $L_J = 0.5$ and $\epsilon_{\min} = 0.05$ rad:
\begin{equation}
\Delta \leq \sqrt{\frac{2 \times 0.05}{0.5}} \approx 0.45 \text{ rad}
\end{equation}
Our default $\Delta = 0.2$ rad satisfies this with a safety margin.

\textbf{Convergence Tolerance:}
Theorem 1 predicts linear convergence, but the rate depends on problem conditioning. In practice, we observe:
\begin{itemize}
\item Fast initial decrease (first 5--10 iterations)
\item Slow asymptotic convergence (iterations 10--50)
\item Diminishing returns after $\|\delta\|_\infty < 10^{-4}$
\end{itemize}

This motivates using relaxed tolerances ($10^{-4}$) for real-time applications.

\textbf{Constraint Satisfaction:}
The $O(\Delta^2)$ mismatch bound explains the small tube excess observed in experiments (typically $< 10^{-3}$ rad with $\Delta = 0.2$).

\subsection{Complexity Analysis}

\begin{proposition}[Per-Iteration Complexity]
Each outer iteration of the sequential convexification algorithm has complexity:
\begin{equation}
O(M \cdot \text{cost}_{\text{ADMM}})
\end{equation}
where the inner ADMM solver requires:
\begin{itemize}
\item One sparse linear system solve: $O(M)$ (banded structure)
\item Per-sample projections: $O(M)$ (parallelizable)
\item Typically 50--200 iterations to convergence
\end{itemize}

Total complexity: $O(K \cdot I \cdot M)$ where $K$ is outer iterations and $I$ is average ADMM iterations.
\end{proposition}

In practice, $K \approx 15$--$30$ and $I \approx 50$--$100$, giving empirical complexity approximately linear in $M$ for the problem sizes tested ($M \leq 10^4$).
